% from main.tex
%%%%%%%%%%  scan idx 423  |  folio 410  %%%%%%%%%%
\textbf{Corollaire 6.3.} Avec les hypothèses de 6.1, soit $A$ un anneau commutatif. Pour qu'un élément $y$ de $K^{\bullet}(X)\otimes_{\mathbb Z}A=K^{\bullet}(X)_A$ soit inversible, il faut et il suffit que $z=(\varepsilon\otimes\Id_A)(y)$ le soit. En particulier, si $E$ est un faisceau localement libre de type fini sur $X$, partout $\ne0$, alors sa classe dans $K^{\bullet}(X)_{\mathbb Q}$ est inversible.

L'idéal d'augmentation de $K^{\bullet}(X)\otimes_{\mathbb Z}A$ est $\Fil^1K^{\bullet}(X)\otimes_{\mathbb Z}A$, et est donc un nilidéal, d'après 6.1. Pour qu'un élément de $K^{\bullet}(X)\otimes_{\mathbb Z}A$ soit inversible il est donc nécessaire et suffisant que son augmentation le soit.

\textbf{Corollaire 6.4.} Soient $X$ un schéma quasi-compact, ayant un faisceau ample, et $f:X'\to X$ un morphisme propre et parfait. On suppose de plus vérifiée l'une des deux conditions suivantes:
\[
  \text{i) } f\text{ projectif},\qquad
  \text{ii) } X\text{ noethérien},
\]
de façon que le «théorème de finitude» soit vrai (III 2.2). On se donne un anneau $A$, et on considère l'homomorphisme «image directe» (IV 2.11):
\[
  f_*:K^{\bullet}(X')\otimes_{\mathbb Z}A\longrightarrow K^{\bullet}(X)\otimes_{\mathbb Z}A.
\]
On suppose enfin qu'il existe un élément $x'\in K^{\bullet}(X')\otimes_{\mathbb Z}A$ tel que $\varepsilon(f_*(x'))\in H^0(X,A)$ soit inversible. Alors l'homomorphisme :
\[
  f^*:K^{\bullet}(X)\otimes_{\mathbb Z}A\longrightarrow K^{\bullet}(X')\otimes_{\mathbb Z}A
\]
est injectif.

On écrit la formule de projection:
\[
  f_*(f^*(x)x')=x.f_*(x').
\]
Or $f_*(x')$ est d'augmentation inversible, donc est inversible (6.3). Par suite, l'homomorphisme composé :
